% LSPV-4 REPORT::MODELBRIDGE-0024::2026-10-12T11:55:57Z
\documentclass[11pt]{article}
\usepackage[margin=1in]{geometry}
\usepackage{amsmath,amssymb,booktabs,array,enumitem,hyperref}
\hypersetup{colorlinks=true,linkcolor=blue,urlcolor=blue}
\newcommand{\Pclass}{\mathsf P}
\newcommand{\NPclass}{\mathsf{NP}}
\newcommand{\OPEN}{\textbf{OPEN}}
\newcommand{\EXACT}{\textbf{PROVED/EXACT}}
\newcommand{\KNOWN}{\textbf{KNOWN THEOREM}}
\newcommand{\COMP}{\textbf{COMPUTATIONAL}}
\newcommand{\RETIRED}{\textbf{FALSIFIED/RETIRED}}
\title{MODELBRIDGE-0024: A polynomial complexity core on a fixed clean core is exactly separation strength}
\author{LSPV-4 / UNIVERSAL\_CORE\_BRIDGE}
\date{12 October 2026, 11:55 UTC}

\begin{document}
\maketitle

\section*{Executive result}
\textbf{P versus NP remains OPEN.}  This run tests whether an infinite
recursive ``clean-core'' set on which every correct uniform decider is
eventually slower than every polynomial can be obtained as a
solver-independent intermediate.  It cannot: for any fixed recursive
NP-complete clean-core language \(K\), existence of such a set is equivalent
to \(K\notin\Pclass\), hence to \(\Pclass\ne\NPclass\).  The equivalence is a
direct corollary of Lynch's 1975 polynomial-complexity-core theorem and
NP-completeness.  Therefore
\texttt{COMPLEXITY-CORE-AS-BRIDGE1} is \RETIRED.  Under the working hostile
hypothesis \(\mathrm{SAT}\in\Pclass\), no polynomial complexity core for
\(K\) exists.

This result sharpens the boundary around LSPV-3's
\texttt{FULL-DOMAIN-SEMANTIC-ANTI-SHARING1}.  If that candidate is
strengthened to one cofinal recursive family on which every decider is
almost-everywhere superpolynomial, it is already the desired separation.  If
it is weakened to incompatibility-table exposure, representation size, or a
named internal resource, the missing \textsc{same\_realization} transfer to
time in an unconditional model remains \OPEN.

\section{Exact language, encoding, and quantifiers}
Fix a self-delimiting binary encoding of signed Monotone
\(3\)SAT-\((2,2)\), and let \(K\subseteq\{0,1\}^*\) be the resulting total
decision language, including a fixed convention for malformed strings.
D\"ocker proves the underlying restricted satisfiability problem NP-complete;
the encoding and malformed-string convention are polynomially interconvertible
with the standard formulation.  Length below is encoded bit length.

A set \(C\subseteq\{0,1\}^*\) is a \emph{polynomial complexity core} for
\(K\) when \(C\) is infinite and recursive and
\[
 \forall D\in\operatorname{Dec}_{\rm det,unif}(K)\;
 \forall p\in\mathbb Z[x]\text{ polynomial}\;
 \exists N\;\forall x\in C\;( |x|\ge N\Rightarrow T_D(x)>p(|x|)).
\]
The set is a subset of the full encoded domain, not silently a proper
yes-instance core.  No advice, oracle, preprocessing table, or uncharged
recognizer is permitted.

For comparison, the final worst-case target previously isolated in COREEQUIV1
has quantifiers
\[
 \forall D\;\forall d\;\forall N_0\;\exists x
 \quad |x|\ge N_0\ \wedge\ T_D(x)>|x|^d.
\]
Lynch's theorem says that nonmembership in \(\Pclass\) permits the stronger
organization of these witnesses into one infinite recursive set \(C\) which
works almost everywhere for every \(D\) and polynomial \(p\).

\section{Clean-core complexity-core equivalence}
\paragraph{Theorem (CLEANCORE-COMPLEXITY-CORE-EQUIVALENCE1).}
For the fixed recursive NP-complete language \(K\), the following are
equivalent:
\begin{enumerate}[leftmargin=*]
\item \(K\) has an infinite recursive polynomial complexity core;
\item \(K\notin\Pclass\);
\item \(\Pclass\ne\NPclass\).
\end{enumerate}

\paragraph{Proof.}
Lynch's Theorem 1 states that every recursive language outside \(\Pclass\)
has an infinite recursive set on which every deterministic decider exceeds
every polynomial almost everywhere.  Thus (2) implies (1).  Conversely, if
\(K\in\Pclass\), choose a total decider \(D_0\) and constants \(c,d\) with
\(T_{D_0}(x)\le c|x|^d+c\) for every sufficiently long \(x\).  A polynomial
larger than that bound defeats the defining condition on every infinite
candidate \(C\), so (1) implies (2).  Finally, \(K\) is in \(\NPclass\) and is
NP-complete under polynomial-time many-one reductions, hence
\(K\notin\Pclass\) iff \(\Pclass\ne\NPclass\). \hfill\(\square\)

\paragraph{Proof status.}
Lynch's result is a \KNOWN{} theorem.  The displayed specialization and the
two elementary converse implications are \EXACT{} pending independent review.
It is not a proof of either side of P versus NP because no unconditional proof
of item (1) or item (2) is supplied.

\section{Semantic compatibility and the failed bridge}
For a CNF \(F\) and consistent partial assignment \(\alpha\), define
\[
 H(F,\alpha)=1\quad\Longleftrightarrow\quad
 \exists\beta\supseteq\alpha\;[\beta\models F].
\]
The full function \(H_F\) can have an exponentially large truth table.  That
semantic-state count does not imply runtime, memory, communication, or
materialization by an arbitrary one-bit decider.  If SAT is in P, any selected
value \(H(F,\alpha)\) is computable in polynomial time by restricting \(F\)
and calling the assumed decider; a witness can likewise be recovered by
self-reduction.  Both are new charged executions, not evidence about the
internal realization of the original run \(D(F)\).

Suppose one tried to repair exposure claims by asserting a recursive clean
family \(C\) on which \emph{every} decider must pay more than every
polynomial.  The theorem shows that this repair is not a weaker transfer:
it is equivalent to the separation.  Under \(\mathrm{SAT}\in\Pclass\), the
restricted language \(K\) has a global polynomial decider, so the proposed
family cannot exist.  This is an exact hostile-model counterexample, not an
inference from easy islands.

\section{Hostile controls and finite regression audit}
The deterministic script
\path{experiments/review/complexity_core_equivalence_synthesis.py} audits the
finite De Morgan analogue of the quantifier pattern for lengths \(2\) through
\(12\), exponents \(1\) through \(6\), seven cost schedules, and every
nonempty finite length subset.  It separately checks fourteen families
(Horn, 2SAT, affine/Tseitin, bounded treewidth, parity, signed clean-core toy
instances, padding, redundant clauses, independent products, alternate
encodings, SAT, UNSAT, duplicate clauses, and mixed controls) under DIRECT,
COPY, DEAD, DELAY, INTERP, and RECOMP wrappers.  It also checks a finite
singleton recognizer to prevent an easy-island residual from being mistaken
for a uniform decider.

The run completed in \(0.587\) seconds without chunking and checked
\(101{,}573\) assertions, \(12{,}544\) singleton-recognizer comparisons, and
\(672\) wrapper cases, with zero failures.  Script SHA-256 is
\texttt{4cf41812440ded65b1d054d12f86c81fa3e07c678f07bbf8b22a345d7c9b3fe2};
result SHA-256 is
\texttt{f7db3f642d9ab64032372536d729a0fe76ea9cc8b3eb11f4143b25945a82bfb3}.
The result is at
\path{results/raw/review/2026-10-12_complexity_core_equivalence.json}.
This evidence is \COMP{} only; it neither proves Lynch's theorem nor
extrapolates a lower bound.

\section{Primary-literature and barrier-transfer preflight}
The two blocking obstacles are (i) obtaining an unconditional complexity core
without assuming \(K\notin\Pclass\), and (ii) transferring semantic structure
to a fully charged resource of the same arbitrary computation.
\begin{itemize}[leftmargin=*]
\item Lynch's polynomial complexity-core theorem is a
\textbf{directly transferable theorem}: its objects and deterministic-time
model match the proposed cofinal clean-family obligation exactly.
\item D\"ocker's NP-completeness theorem for Monotone \(3\)SAT-\((2,2)\) is
directly transferable after the signed encoding is fixed.  It removes any
need for an additional transfer from the fixed clean core to unrestricted
formula families.
\item Baker--Gill--Solovay relativization and Aaronson--Wigderson
algebrization remain known barriers.  The elementary equivalence and the
semantic restriction query relativize; no nonrelativizing separation step is
present.
\item FLMV time--space lower bounds transfer only after one ordinary machine
is proved to occupy their exact simultaneous resource region.  A complexity
core alone gives no such fixed landing.
\item Feasible interpolation, VLSI congestion, control observability,
communication cuts, coding-theoretic local decoding, and local-to-global
gluing are useful analogies only if a new interface-preserving theorem maps
the arbitrary decider to their restricted object.  No such theorem is known
here.  Treating similar vocabulary as a model transfer is rejected.
\end{itemize}

\section{Barrier and wrapper audit}
\paragraph{SAME\_REALIZATION.}
Lynch's theorem is already a statement about the time of each complete
decider; it needs no internal-state extraction.  In contrast, an \(H\)-table,
frontier, protocol transcript, proof node, branching program, circuit, or
Gaussian-elimination state is not forced into that realization.

\paragraph{Wrappers and alternate representations.}
COPY, DEAD, DELAY, INTERP, and RECOMP preserve the language while changing
traces and named intermediates.  They do not evade a true end-to-end time
core, but they refute any claim that the core forces one particular internal
object.  Horn, affine, 2SAT, and bounded-treewidth algorithms likewise show
why semantic size is not a cost measure.

\paragraph{Arbitrary polynomial exponents and nonuniformity.}
The definition quantifies over every polynomial, so it does not hide a fixed
exponent.  Conversely, choosing a different hardwired decider or table for
each length introduces advice and is not a uniform counterexample.  Padding
cannot create a separation: encoded length and all padding work remain
charged.

\paragraph{Natural Proofs and proof systems.}
No large constructive circuit property, proof-system lower bound, or
circuit-to-algorithm simulation is established.  Razborov--Rudich is neither
invoked nor escaped.  HEC, VSC-HEC, G-PPC-HEC, CRS-PPC, RCCN, and CCSPC retain
their exact restricted-model results but still lack a total compiler from
every static decider and a model-matched unconditional endpoint.

\section{Post-quantum technique ledger}
\begin{itemize}[leftmargin=*]
\item Schur squares, shortening/puncturing/duals, quadratic relation ranks,
Hilbert functions, syzygies, and Betti numbers: \textbf{FAILED TRANSFER} for
arbitrary one-bit deciders.  Such structure is not forced and may be an
algorithmic weakness.
\item Collision entropy: \textbf{EXACT OBJECT/THEOREM TRANSFER} only for a
declared map \(L\), where
\( |\operatorname{im}L|\ge |A|^2/\operatorname{Coll}(L) \).  It gives no
end-to-end time core and no original-run exposure.
\item LWE-style worst-case/average-case extraction:
\textbf{ANALOGY/CANDIDATE needing a new bridge}.  A black-box extractor is a
new realization and known NP worst-case/average-case barriers must be met.
\item Collision extraction/special soundness:
\textbf{ANALOGY/CANDIDATE needing a construction-specific relation}; generic
cryptographic binding is inadmissible.
\item Hash/random-oracle methods: \textbf{FAILED TRANSFER} unless all
hardness and random-oracle assumptions are removed.
\end{itemize}

\section{CORPUS REVIEWED}
The run enumerated the current GitHub tree and read root \path{main.tex},
actual \path{main4.tex} and \path{main5.tex}, \path{AGENTS.md},
\path{CODEX.md}, the pnp-research skill and protocol references, and every
current task specification.  It inspected verified MAIN SYNTHESIS-0023, the
latest LSPV-1 CIRCEQUIV1, LSPV-2 REPDIRSUM1, predecessor LSPV-3 ANTISHARE1,
MODELBRIDGE-0023, materially relevant COREEQUIV1, RESOURCEAXIS1, QGUARD1,
COVCOLL1, CALQNT1 and earlier universal-bridge reports, and the
HEC/VSC-HEC/G-PPC-HEC/CRS-PPC/CCSPC causal branch.  Governance notes and
cited experiment/result artifacts were surveyed, and repository search was
used to locate duplicates, contradictions, superseded assumptions, failed
transfers, and retired routes.

The immediately preceding LSPV-3 ANTISHARE1 report exists in its canonical
directory, is substantive, contains the required header and sentinel, and
was read back from GitHub.  No recovery request was needed.  The exact body
of MODELBRIDGE-0022 also already exists and verifies at the required canonical
LSPV-4 path; it was not duplicated or reconstructed.

\section{DUPLICATION / CONTRADICTION AUDIT}
COREEQUIV1 already proved that an all-decider, all-polynomial worst-case lower
bound on an NP-hard encoded family is equivalent to \(P\ne NP\).  ANTISHARE1
then isolated a full-domain semantic anti-sharing obligation.  The new
increment is the primary-literature identification of Lynch's stronger
uniformization: the cofinal witnesses can be collected into one infinite
recursive almost-everywhere hard set exactly when the language is outside P.
Thus the apparent ``clean complexity core'' repair is not a weaker missing
bridge.  No verified result is contradicted.

\section{WHAT CHANGED RELATIVE TO ROOT MAIN}
Root MAIN already records the universal-core transfer and model-matched
landing as \OPEN.  This report adds the exact complexity-core classification,
retires complexity-core existence as an intermediate, and narrows the
surviving target to either a nonrelativizing direct separation or a genuine
same-realization transfer into a restricted unconditional endpoint.  Actual
\path{main4.tex} and \path{main5.tex} remain promise/partial-domain ARC and
WRAPPER construction manuscripts; they are not silently redefined as total
arbitrary-algorithm compilers.

\section{ADOPTED / REJECTED / RETIRED PRIOR COMPONENTS}
\begin{itemize}[leftmargin=*]
\item \textbf{ADOPTED:} corrected signed NP-complete clean core; encoded
length; semantic \(H(F,\alpha)\); Lynch complexity cores; full charging;
hostile wrappers; \textsc{same\_realization}; absolute model matching.
\item \textbf{REJECTED:} semantic-state count as runtime; proper-core status
without proof; named internal APIs; free extraction; padding dilution;
nonuniform tables; cryptographic assumptions.
\item \textbf{FALSIFIED/RETIRED:}
\texttt{COMPLEXITY-CORE-AS-BRIDGE1}.  Its conclusion is equivalent to the
desired separation for the fixed NP-complete language.
\item \textbf{OPEN:} an unconditional proof that this complexity core exists,
or an alternative noncircular arbitrary-decider transfer with a
model-matched forbidden landing.
\end{itemize}

\section{STALE OR OVERSTATED CLAIMS FOUND}
Any statement that a fixed clean family ``must'' be a polynomial complexity
core before proving \(P\ne NP\) is circular unless it supplies a genuinely
new proof of the separation.  Any claim that Lynch's set lies entirely in the
yes-instances is also overstated; the theorem used here only requires a
recursive subset of the full encoded domain.  Any inference from the size of
\(H_F\) to decider time, or from externally queryable \(H\) values to
original-run exposure, remains invalid.

\section{Strict status ledger and handoff}
\begin{center}
\begin{tabular}{p{0.49\linewidth}p{0.20\linewidth}p{0.23\linewidth}}
\toprule
Item & Status & Reason \\
\midrule
Lynch polynomial complexity-core theorem & \KNOWN & Primary theorem, 1975\\
\texttt{CLEANCORE-COMPLEXITY-CORE-EQUIVALENCE1} & \EXACT & Lynch plus NP-completeness\\
Finite quantifier/wrapper audit & \COMP & 101,573 assertions; zero failures\\
\texttt{COMPLEXITY-CORE-AS-BRIDGE1} & \RETIRED & Equivalent to \(P\ne NP\)\\
Semantic-object to arbitrary-run resource transfer & \OPEN & No forced representation\\
P versus NP & \OPEN & No unconditional item (1) or (2)\\
\bottomrule
\end{tabular}
\end{center}

\paragraph{Precise surviving target.}
Do not ask an arbitrary decider to expose \(H\) or posit a complexity core as
an intermediate.  Either prove directly, without assuming the separation,
that the fixed NP-complete language lies outside P, or give a total uniform
source-defined compiler that maps every hypothesized polynomial clean-core
decider to the same fully charged realization in a cited unconditional
forbidden model.  The compiler must survive alternate encodings, arbitrary
polynomial exponents, all hostile wrappers, and relativization/algebrization
audits.  This target is \OPEN.

\begin{thebibliography}{9}
\bibitem{Lynch75} N. A. Lynch, ``On Reducibility to Complex or Sparse Sets,''
\emph{Journal of the ACM} 22(3), 341--345, 1975.
\bibitem{Docker20} J. D\"ocker, ``Monotone 3-Sat-(2,2) is NP-complete,''
\emph{Theoretical Computer Science} 802, 1--9, 2020.
\bibitem{Cook71} S. A. Cook, ``The complexity of theorem-proving
procedures,'' STOC 1971, 151--158.
\bibitem{BGS75} T. Baker, J. Gill, and R. Solovay,
``Relativizations of the P=?NP question,'' \emph{SIAM J. Comput.} 4(4),
431--442, 1975.
\bibitem{RR97} A. Razborov and S. Rudich, ``Natural Proofs,''
\emph{J. Comput. Syst. Sci.} 55(1), 24--35, 1997.
\bibitem{AW09} S. Aaronson and A. Wigderson, ``Algebrization: A New Barrier
in Complexity Theory,'' \emph{ACM TOCT} 1(1), 2009.
\bibitem{FLMV05} L. Fortnow, R. Lipton, D. van Melkebeek, and A. Viglas,
``Time-space lower bounds for satisfiability,'' \emph{J. ACM} 52(6),
835--865, 2005.
\end{thebibliography}

\noindent\textbf{Final verification sentinel:} report body complete; exact
quantifiers and all charges stated; corpus and predecessor checked; strict
statuses preserved; P versus NP remains \OPEN.
% END_OF_VERIFIED_REPORT::LSPV-4::MODELBRIDGE-0024
\end{document}
